% !TeX root = ../main.tex
\section{引言}
电子结构计算为分子结构、光谱与反应性研究提供了重要工具。密度泛函理论（density functional theory, DFT）通过电子密度描述基态多电子体系，Kohn--Sham 方案使其能够转化为实际可求解的自洽问题\cite{hk,ks}。然而，计算软件能够完成某种求解，并不意味着研究者的科学问题已经得到回答。方法和基组选择、电荷与自旋设定、结构准备、收敛诊断以及结果解释共同决定计算的有效性。

传统自动化主要将已知计算协议转化为可重复执行的流程。大语言模型则提供了进一步连接自然语言需求、软件知识与程序生成的可能性。ReAct、Toolformer 与 Reflexion 分别从推理与行动交替、工具使用及反馈修正等角度发展了语言模型 Agent 的方法基础\cite{react,toolformer,reflexion}。ChemCrow 和 Coscientist 展示了模型调用化学工具、组织化学研究操作的可行性\cite{chemcrow,coscientist}。对于量子化学，这一方向的关键问题已从“能否生成输入”扩展为“能否保持科学目标，并根据中间证据组织后续计算”。

本文沿 BG6022 V2、V3、V4 的顺序考察三种研究思路。V2 聚焦计算链的可追溯性，首先建立受控执行与证据复核的基础。V3 将研究对象扩展到有限任务范围内的自然语言交互和工具组合。V4 则将任务特定算法的生成置于中心，研究如何在统一工具边界内，以生成代码处理更开放的科学目标。该顺序表示研究重心的变化，不预设后续版本在所有任务上具有更高成功率。

V4 是本文的主要研究对象。系统采用 ORCA 及 ORCA Python Interface（OPI）作为科学计算基础\cite{orca2020,orca6,opipaper}，使用单一决策循环组合通用工具与现场代码，并将原始目标、实际执行、科学检查和结果交付分别记录。本文的贡献包括：建立三个版本之间的研究问题联系；阐述 V4 的目标保持与受控代码执行机制；通过科学案例和跨任务评测分析执行成功与完整目标完成之间的差异。本文关注自主计算系统的方法与验证，不提出新的电子结构近似。
